\clearpage
\section{Continuity with the 1962 Missal and Postconciliar Developments}
\subsection*{Summary of Continuity and Change from the 1962 Roman Missal}
{\footnotesize\begin{longtable}{>{\raggedright\arraybackslash}p{0.21\linewidth}>{\raggedright\arraybackslash}p{0.19\linewidth}>{\raggedright\arraybackslash}p{0.53\linewidth}}
\toprule
\textbf{Object} & \textbf{Relationship} & \textbf{Direct edition comparison} \\
\midrule
\endhead
\textbf{Altar} & \PCDTag{materially changed} & Whereas the 1962 Roman Missal ordinarily presupposed a stone altar (or recessed consecrated altar stone holding relics inside the mensa) dressed with three cloths, candlesticks, central crucifix, three altar cards, and normally the tabernacle, the 2008/2011 Missal recommends a single freestanding altar wherever possible (GIRM 299, 303), deposits saints' relics beneath the altar (GIRM 302), clears the mensa of altar cards and extraneous items (GIRM 306), and in the United States permits solid wood for a fixed altar mensa (US GIRM 301). \\
\textbf{Ambo} & \PCDTag{newly provided} & The 1962 Roman Missal did not prescribe a permanent sanctuary ambo: at Low Mass the priest read the Epistle and Gospel from opposite corners of the altar, and at Solemn Mass the subdeacon and deacon sang them from the sanctuary floor with the subdeacon holding the Gospel book (RS VI.4--5). Postconciliar law restores a dedicated stationary ambo as a standard architectural appointment of every sanctuary (Inter Oecumenici 96; GIRM 309). \\
\textbf{Presidential Chair} & \PCDTag{materially changed} & In the 1962 Roman Rite, presbyteral seating consisted of a three-seat bench (sedilia) placed along the Epistle side of the sanctuary solely for sitting during sung chants and lessons while all orations were spoken at the altar. Postconciliar law replaces those lateral seats with a focal single chair from which the priest leads the Introductory Rites, Liturgy of the Word, and Concluding Rites (Inter Oecumenici 92; GIRM 310). \\
\textbf{Credence Table} & \PCDTag{continued} & Both the 1962 Roman Missal (RG 528; RS II.5) and the 2008/2011 GIRM (118c, 139, 279) appoint a sanctuary side table (abacus) for cruets, lavabo articles, Communion-plate, and vessels. In postconciliar Mass with a congregation, its role expands because the chalice and linens also ordinarily wait there until the Preparation of the Gifts (rather than being carried to the altar by the priest at the start of Low Mass) and vessel purification commonly takes place at the credence table (GIRM 279). \\
\textbf{Tabernacle} & \PCDTag{materially changed} & Both the 1917 Code (cann. 1268--1269) and current law (CIC 1983 can. 938; GIRM 314) mandate a single immovable, solid, opaque, securely locked tabernacle. Whereas 1917/1962 discipline ordinarily fixed the veiled tabernacle in the middle of the high altar of celebration, GIRM 315 considers it clearer as a sign not to place it on the altar where Mass is celebrated, entrusting the Diocesan Bishop to choose either another worthy place in the sanctuary or an organically connected adoration chapel, and no longer universally mandating a tabernacle veil. \\
\textbf{Sanctuary Lamp} & \PCDTag{continued} & Postconciliar discipline continues the 1917/1962 requirement without interruption: both CIC (1917) can. 1271 and CIC (1983) can. 940 together with GIRM 316 require a special living flame of oil or wax to burn continually before or near the tabernacle whenever the Eucharist is reserved. \\
\textbf{Baptismal Font} & \PCDTag{materially changed} & Both 1962 discipline and CIC (1983) can. 858 require a baptismal font in every parish church. Whereas 1962 practice ordinarily enclosed the font inside a gated baptistery near the church entrance and blessed the year's water at the Easter Vigil with an infusion of Oil of Catechumens and Sacred Chrism directly into the basin, postconciliar rites bless fresh water at each Baptism outside Easter Time without pouring oils into the font and favor congregationally visible placement suited to pouring or immersion. \\
\textbf{Chalice} & \PCDTag{materially changed} & Whereas the 1962 Roman Missal required a gold or silver-cupped chalice consecrated with Sacred Chrism by a bishop and carried into Low Mass veiled with a stiff burse by the priest himself (RS I.1, II.1), the 2008/2011 Missal appoints a liturgical blessing by bishop or priest (GIRM 333), ordinarily prepares the chalice on the credence table until the Preparation of the Gifts (GIRM 118c, 139), and commends the chalice veil as a praiseworthy optional practice while leaving the burse unnamed. \\
\textbf{Paten} & \PCDTag{materially changed} & In the 1962 Roman Missal, the paten was a small flat gilded metal disc consecrated by a bishop, sized for the celebrant's single host, kissed and used to sign the celebrant at the Libera nos, and held veiled by the subdeacon at Solemn Mass (RS I.1, VII.9, X.1--2). The 2008/2011 GIRM omits the signing and kissing at the Libera nos, introduces a broad paten fittingly holding hosts for clergy and faithful alike (GIRM 331), and under USCCB GIRM 329 admits solid noble non-metal materials. \\
\textbf{Ciborium} & \PCDTag{continued} & Both the 1962 Roman Missal (using Latin pyxis or vas) and the 2008/2011 IGMR (continuing Latin pyxis, rendered ``ciborium'' in English) appoint a blessed covered vessel for consecrating, distributing, and reserving the Sacred Hosts. Postconciliar norms expand distribution use to broad dignified bowls alongside stemmed ciboria, admit approved non-metal solid noble materials in the United States (US GIRM 329), and treat the silk ciborium veil as customary rather than universally mandated. \\
\textbf{Pyx} & \PCDTag{continued} & The portable hinged metal pyx carried in a pouch for Communion of the sick continues unchanged in construction from the 1962 Roman Ritual (Title IV, ch. 2) into postconciliar discipline (GIRM 329; CIC can. 939). Postconciliar pastoral law narrows English usage of ``pyx'' specifically to this portable sick-call container and extends its authorized pastoral carriage to formally deputed extraordinary ministers alongside clergy. \\
\textbf{Communion Cup} & \PCDTag{newly provided} & In the 1962 Roman Missal, lay faithful received Holy Communion exclusively under the species of bread (RG 503; RS X.6), so no supplementary Communion chalices existed. Following Sacrosanctum Concilium 55, GIRM 281--287 and the USCCB Norms newly provide for additional chalices governed by the full sacred-vessel standards of GIRM 328--330. \\
\textbf{Cruets} & \PCDTag{continued} & Both the 1962 Roman Missal (using Latin ampullae, RG 528) and the 2008/2011 IGMR (using Latin urceoli, GIRM 118c) appoint paired wine and water cruets on the credence table for the Offertory mixing and ablutions. Postconciliar rubrics omit kissing the cruets and signing the water cruet with the cross, and allow the wine to be carried from the nave in the Procession of the Gifts (GIRM 140). \\
\textbf{Flagon} & \PCDTag{no established counterpart} & The 1962 Roman Missal had no counterpart to a wine flagon because lay Communion was distributed under the species of bread alone and small cruets on the credence table sufficed (RG 528). The flagon emerged with the postconciliar gifts procession and Communion under both kinds, and was strictly restricted to pre-consecration wine transport by Redemptionis Sacramentum 106 (2004). \\
\textbf{Altar Cloth} & \PCDTag{materially changed} & The 1962 Roman Missal required three clean blessed white linen cloths covering the altar, with the uppermost cloth reaching the ground at both sides (RG 526). Postconciliar law requires at least one white cloth suited to the altar's shape and size (GIRM 117, 304), while the United States adaptation permits additional festive colored under-cloths beneath the white top cloth. \\
\textbf{Corporal} & \PCDTag{continued} & The square white linen corporal folded inward in nine squares continues directly from the 1962 Roman Missal (RS I.1). In postconciliar Mass with a congregation, it is placed on the altar at the Preparation of the Gifts rather than drawn from a stiff burse at the start of Low Mass (1962 RS II.2), and GIRM 118c no longer mandates a burse. \\
\textbf{Purificator} & \PCDTag{continued} & The folded white linen purificator continues directly from the 1962 Roman Missal (RS I.1, X.5), where a single purificator served the celebrant's chalice and ablutions. In postconciliar celebration, Communion under both kinds expands its use so that a separate clean purificator accompanies every Communion chalice (GIRM 286). \\
\textbf{Pall} & \PCDTag{optional} & The 1962 Roman Missal required a blessed small linen pall on every prepared chalice stack to cover the host and chalice throughout the Offertory and Canon (RS I.1, VII.5, VIII.7). Postconciliar law expressly makes the pall optional, prescribing it ``if appropriate'' (pro opportunitate, GIRM 118c). \\
\textbf{Lavabo Towel} & \PCDTag{continued} & The hand towel (manutergium) prepared with water and basin on the credence table continues directly from the 1962 Roman Missal (RG 528; RS VII.6--7, VII.10). In the postconciliar rite, the priest washes his hands rather than only thumb and forefinger tips, and replaces Psalm 25:6--12 (Lavabo inter innocentes) with the single verse of Psalm 50:4 [51:2] praying for interior cleansing (GIRM 76, 145). \\
\textbf{Roman Missal} & \PCDTag{materially changed} & The 1962 Missale Romanum was a plenary altar book holding all presidential orations, chants, Epistles, and Gospels, remaining on the altar throughout Mass and carried between the Epistle and Gospel corners (RG 527; RS VI.5, XI.1). In the postconciliar Roman Rite, the Scripture readings are separated into the Lectionary and Book of the Gospels, while the Missal serves at the presidential chair and rests on the altar only from the Preparation of the Gifts through purification (GIRM 118a, 139, 306). \\
\textbf{Lectionary for Mass} & \PCDTag{materially changed} & In the 1962 Roman Missal, the one-year cycle of Epistles and Gospels was printed inside the Missal itself, with an optional separate Epistle book (liber Epistolarum) for the subdeacon at Solemn Mass (RS VI.4). The postconciliar reform established a dedicated multi-volume Lectionary spanning a three-year Sunday cycle (A, B, C) and two-year weekday cycle (I, II), placed open on the ambo for every Mass with a congregation (GIRM 118b; OLM 1981). \\
\textbf{Book of the Gospels} & \PCDTag{changed use} & In the 1962 Roman Missal, a separate Gospel book (liber Evangeliorum) belonged to Solemn Mass, where it was placed on the altar only before the Gospel, carried by the deacon in the Gospel procession, held by the subdeacon between lighted candles, incensed, and brought to the celebrant to kiss (RS VI.5). Postconciliar rubrics integrate the Book of the Gospels into ordinary Sunday and festive celebrations, carrying it slightly elevated in the entrance procession and resting it on the altar from the start of Mass until the Gospel procession (GIRM 120d, 122, 133--134, 306). \\
\textbf{Book of the Chair} & \PCDTag{no established counterpart} & Because the priest in the 1962 Roman Missal prayed the Collect and Postcommunion at the altar rather than at a presidential chair (RG 527), there was no presbyteral Book of the Chair (only a minister holding the Missal for a bishop under Caeremoniale Episcoporum I.20). The Holy-See-confirmed USCCB volume Excerpts from the Roman Missal (2018) is a distinct postconciliar practical companion fulfilling GIRM 118a. \\
\textbf{Missal Stand} & \PCDTag{continued} & The Missal stand or cushion continues in practical form from the 1962 Roman Missal, which expressly required a cushion or bookstand beneath the Missal on the altar from the start of Mass and directed the server to move it from the Epistle side to the Gospel side and back (RG 527; RS VI.5, XI.1). In the postconciliar Roman Rite, it is neither mandated by name nor kept on the altar before the Preparation of the Gifts, but set once on the mensa for the Liturgy of the Eucharist and removed afterward (GIRM 139, 306). \\
\textbf{Processional Cross} & \PCDTag{changed use} & Whereas the 1962 Roman Missal required a separate crucifix standing in the center of the altar (RG 527) and did not ordinarily prescribe a carried cross at a priest's High Mass entrance, the 2002/2008 GIRM assigns the processional cross with the figure of Christ crucified a regular place in the entrance procession (GIRM 120b) and permits it to stand near the altar as the sole altar cross (GIRM 122, 308). The 2002 third typical edition also restored the explicit requirement of the sculpted corpus that had been omitted in the 1970 and 1975 editions. \\
\textbf{Altar Candles} & \PCDTag{materially changed} & Both the 1962 Roman Missal (RG 527) and the 2002/2008 GIRM (117, 307) scale candle counts from two up to six, or seven for the Diocesan Bishop. In 1962, altar candlesticks stood directly on the altar or its gradine flanking the central crucifix, supplemented by separate Sanctus and elevation torches; postconciliar GIRM 117, 122, and 307 permit candlesticks to stand either on or around the altar---including processional candlesticks placed near it---so that the faithful maintain an unobstructed view of the altar table. \\
\textbf{Sanctuary Bells} & \PCDTag{optional} & Whereas the 1962 Roman Missal required a small bell on the Epistle-side table and prescribed ringing at the Sanctus and elevations (RG 528; RS VII.8, VIII.6), the 2002/2008 GIRM makes the bell optional according to local custom---shortly before the consecration and at each showing of the Host and Chalice---omitting the mandatory Sanctus and Communion rings (GIRM 150). \\
\textbf{Aspergillum and Aspersorium} & \PCDTag{changed use} & Although the portable metal bucket and sprinkler continue unchanged in form, the 1962 Roman Missal placed the Sunday Asperges me or Vidi aquam before the start of Mass, with the water blessed in the sacristy and the celebrant wearing a cope before putting on the chasuble. In the 2002/2008 Roman Missal, the Blessing and Sprinkling of Water takes place inside the Introductory Rites in place of the Penitential Act, celebrated by the priest already vested in his chasuble (GIRM 51, 118c; Appendix II). \\
\textbf{Thurible} & \PCDTag{changed use} & In the 1962 Roman Missal, incense at Mass was restricted to Solemn Mass and permitted sung Masses (never Low Mass) and accompanied by spoken blessing and incensation prayers (Ordo Missae 1035--1038). The 2002/2008 GIRM permits incense optionally in any form of Mass at five enumerated moments, replaces the spoken blessing with a silent Sign of the Cross, and codifies standardized sets of three, two, and single swings (GIRM 276--277). \\
\textbf{Incense Boat} & \PCDTag{continued} & Both the 1962 Roman Missal (RS IV.4, VII.10) and the 2002/2008 GIRM (120, 276--277) pair the incense boat and spoon with the thurible. In postconciliar practice, the priest imposes incense silently with a Sign of the Cross rather than speaking the 1962 blessing formula, and neither the spoon nor the celebrant's hand is kissed at the presentation. \\
\textbf{Alb} & \PCDTag{changed use} & In the 1962 Roman Missal, the alb was worn solely by ordained sacred ministers at Mass (bishop, priest, deacon, and subdeacon, RS I.2--3), while altar servers and minor clerics wore cassock and surplice (RS II.1, VI.8). Under the 2002/2008 GIRM, the alb becomes the universal common garment of all ordained and instituted ministers, may never be replaced by a surplice under a chasuble, dalmatic, or stole alone (GIRM 336), and in the United States is also permitted for altar servers, readers, and other lay ministers (US GIRM 339). \\
\textbf{Cincture} & \PCDTag{optional} & In the 1962 Roman Missal, the cincture was unconditionally mandatory with every alb both to gather the tunic and to anchor the priest's crossed stole across the breast (RS I.3). Under the 2002/2008 GIRM, the cincture becomes conditionally optional whenever an alb is tailored to fit cleanly without one (GIRM 336), and the priest's stole hangs straight rather than crossing through the cincture (GIRM 340). \\
\textbf{Amice} & \PCDTag{optional} & In the 1962 Roman Missal, the amice was unconditionally required as the first Mass vestment for every bishop, priest, deacon, and subdeacon, kissed at its center cross and placed briefly upon the head before settling on the shoulders (RS I.3). Under the 2002/2008 GIRM, the amice is conditionally required only when the alb fails to cover ordinary attire at the neck (GIRM 336). \\
\textbf{Stole} & \PCDTag{changed use} & While the deacon's diagonal wearing from left shoulder to right hip continues unchanged from 1962, the 1962 Roman Missal required a priest celebrating Mass to cross the stole across his breast and anchor it with the cincture, reserving straight hanging to bishops and prelates (RS I.3--4). Under GIRM 340, every priest now wears the stole hanging straight down in front of his chest. \\
\textbf{Chasuble} & \PCDTag{continued} & Both the 1962 Roman Missal (RG 133--134; RS I.4) and the 2002/2008 GIRM (119a, 337) prescribe the chasuble as the proper Mass vestment of the priest celebrant over alb and stole. Postconciliar rubrics additionally allow the chasuble for rites immediately joined to Mass such as the Sunday sprinkling, Ash Wednesday ashes, and funeral commendation (Tres abhinc annos 26; GIRM 337), emphasize graceful form and noble fabric over heavy ornament (GIRM 344), and introduce a narrow omission faculty for non-principal concelebrants in genuine need (GIRM 209; RS 124). \\
\textbf{Dalmatic} & \PCDTag{continued} & The dalmatic continues from the 1962 Roman Missal (RG 137) as the proper outer vestment of the deacon over alb and diagonal stole. Postconciliar GIRM 338 introduces a limited omission faculty for necessity or lesser solemnity (discouraged by RS 125), while Caeremoniale Episcoporum (1984) no. 56 simplifies episcopal solemn vesting to a single dalmatic beneath the chasuble without a separate tunicle. \\
\textbf{Cassock and Surplice} & \PCDTag{changed use} & In the 1962 Roman Missal, cassock and surplice were the standard vesture of altar servers, choir clergy, and priests administering many non-Mass rites with a stole (RS II.1, VI.8). Postconciliar GIRM 336 replaces surplice-and-stole with alb-and-stole whenever rubrics prescribe a stole alone or beneath chasuble/dalmatic, while retaining cassock and surplice for choir dress (RS 128) and permitting them for altar servers alongside the alb (US GIRM 339; CE 1984 no. 65). \\
\textbf{Cope} & \PCDTag{continued} & The cope continues from 1962 usage in processions, solemn Office, and Eucharistic Benediction (GIRM 341; HCWEOM nn. 91--92). Following Tres abhinc annos (1967) no. 26 and GIRM 51 and 337, postconciliar rubrics permit the chasuble rather than requiring a cope for rites integrated directly into Mass such as the Sunday sprinkling rite and Ash Wednesday imposition of ashes. \\
\textbf{Humeral Veil} & \PCDTag{changed use} & In the 1962 Roman Missal, the humeral veil served both Eucharistic Benediction and processions and---at Solemn High Mass---the subdeacon, who wore it from Offertory to the Pater Noster to hold the empty paten (RS VII.9, X.8). After Paul VI's Ministeria quaedam (1972) suppressed the subdiaconate and the 1969/1970 Order of Mass omitted subdiaconal paten-holding, the humeral veil's liturgical use is dedicated exclusively to Eucharistic Benediction, processions, and Holy Thursday transfer. \\
\textbf{Holy Oils and Oil Stocks} & \PCDTag{materially changed} & The 1962 rites required pure olive oil exclusively, blessed strictly on Holy Thursday morning, and poured both Oil of Catechumens and Sacred Chrism directly into the Easter Vigil baptismal font. Postconciliar discipline (Ordo benedicendi oleum 1970; CIC 1983 can. 847, par. 1) permits oil pressed from olives or other plants, allows moving the Chrism Mass to an earlier day in Holy Week to foster presbyteral participation, discontinues pouring oils into the font water, and encourages prominent reservation in a locked ambry near the baptistery or sanctuary. \\
\bottomrule
\end{longtable}}

\subsection*{Historically Superseded Postconciliar Forms (1967--2011)}
\label{pc-obj-hist-subdeacon-tunicle-maniple}
\PCDHistHead{Subdeacon's Tunicle and Maniple}{Tunicella et manipulus}{1967--1972 --- Tres abhinc annos (1967) no. 25; Ministeria quaedam (1972); 1975 Missale Romanum; CE (1984) no. 56}{N41}
Between 1967 and 1975, two vestments inherited from the preconciliar sanctuary passed out of postconciliar Roman use in stages. First, the Sacred Congregation of Rites Instruction Tres abhinc annos (4 May 1967, effective 29 June 1967), no. 25, declared that the maniple---the ornamental strip of liturgical cloth worn on the left forearm of the priest, deacon, and subdeacon---was no longer required; it ceased to be prescribed and was omitted from the 1970 General Instruction of the Roman Missal rather than being formally banned by decree. Second, the first typical edition of the postconciliar Missale Romanum (26 March 1970, Prot. 166/70) briefly retained the liturgical office of the subdeacon (1970 IGMR nos. 142--152) and his sleeved outer vestment, the tunicle (no. 301). Pope Paul VI's Motu Proprio Ministeria quaedam (15 August 1972, AAS 64:529--534; effective 1 January 1973) then suppressed first tonsure, the four minor orders, and the major order of subdeacon in the Latin Church, establishing the lay instituted ministries of Lector and Acolyte. The second typical edition of the Missale Romanum (27 March 1975) excised the subdeacon and tunicle from the General Instruction, and the Caeremoniale Episcoporum (14 September 1984) no. 56 retained only a single light pontifical dalmatic beneath the bishop's chasuble on solemnities, omitting the separate pontifical tunicle.

\PCDField{Current discipline (2008/2011)}{The 2002/2008 Latin and 2011 United States Roman Missal prescribe neither the subdiaconate nor the tunicle and maniple (GIRM nos. 119, 335--341). Duties formerly assigned to the subdeacon belong to instituted lectors and acolytes vested in the alb (GIRM no. 336; CIC can. 230, par. 1, as amended by Spiritus Domini, 2021), and on major pontifical occasions the bishop wears only a single light dalmatic beneath the chasuble (Caeremoniale Episcoporum 1984 no. 56).}
\PCDField{1962 comparison}{\PCDTag{superseded} In the 1962 Roman Missal, the maniple was required on the left forearm of the priest, deacon, and subdeacon during Mass, and the tunicle served both as the proper outer vestment of the subdeacon at Solemn Mass and as the inner layered garment beneath the bishop's dalmatic and chasuble. Both vestments passed out of postconciliar law between 1967 and 1975 (and from episcopal layering in 1984), remaining in liturgical use only where the 1962 Missal is celebrated.}
\PCDEntryRule

\label{pc-obj-hist-fragile-vessels-pouring}
\PCDHistHead{Fragile Vessels, Flagon Pouring, and EMHC Purification}{Vasa fragilia et infusio}{1970--2006/2011 --- GIRM (2002/2008) nos. 279, 328--330; Redemptionis Sacramentum (2004) nos. 106, 117, 119; CDWDS letter (Oct 2006); USCCB Norms (2011) no. 53}{N42}
The 1970 and 1975 editions of the General Instruction of the Roman Missal (nos. 290--291) omitted precious metal as an explicit primary norm for sacred vessels, directing only that vessels be fashioned from solid materials esteemed as noble in each region (with preference for materials not easily broken or rendered unusable) and that the chalice bowl be nonabsorbent. Over the next three decades, some parishes adopted glass, stoneware, ceramic, or clay chalices and bowls, and consecrated wine in large flagons or pitchers from which the Precious Blood was poured into Communion cups after consecration. The third typical edition (2002/2008 GIRM nos. 328--330) restored precious metal---gilded inside whenever the metal rusts or is less precious than gold---as the universal primary standard, while the United States adaptation (US GIRM no. 329) narrowed secondary solid noble materials such as hard woods to patens, ciboria, and pyxes that do not easily break, deteriorate, or absorb liquids. On 25 March 2004, the Congregation for Divine Worship and the Discipline of the Sacraments issued Redemptionis Sacramentum, which strictly forbade post-consecration pouring of the Precious Blood and barred flagons or bowls from holding the Blood of Christ (no. 106), while explicitly reprobating glass, earthenware, clay, and fragile or rusting vessels (no. 117). Finally, a temporary three-year United States indult (2002--2005) allowing extraordinary ministers of Holy Communion to purify sacred vessels expired without renewal per the Congregation's letter of October 2006, and the 2011 USCCB Norms (no. 53) strictly reserve purification to the priest, deacon, or instituted acolyte.

\PCDField{Current discipline (2008/2011)}{Sacred vessels are made primarily of precious metal gilded inside where necessary (GIRM no. 328), with nonabsorbent bowls for all chalices (no. 330) and the United States adaptation limiting solid noble alternatives to hard, non-breaking, non-absorbent patens, ciboria, and pyxes (US GIRM no. 329). Glass, earthenware, and clay vessels and post-consecration pouring from flagons are strictly prohibited (Redemptionis Sacramentum nos. 106, 117). Wine brought in a flagon is poured into chalices before the Eucharistic Prayer, and vessel purification belongs exclusively to the priest, deacon, or instituted acolyte (GIRM no. 279; USCCB Norms no. 53).}
\PCDField{1962 comparison}{\PCDTag{no established counterpart} Fragile or porous sacred vessels, post-consecration pouring from flagons, and lay vessel purification had no counterpart in the 1962 Roman Missal, which mandated gold or interior-gilded silver chalices and patens, distributed Communion to the laity under bread alone, and restricted ablutions exclusively to the priest celebrant. The 2002/2008 Missal, Redemptionis Sacramentum (2004), and the 2011 United States norms restore precious metal as the universal primary norm, bar fragile materials and post-consecration pouring, and restrict purification to the priest, deacon, or instituted acolyte.}
\PCDEntryRule

\label{pc-obj-hist-1970-sanctuary-clarifications}
\PCDHistHead{Sanctuary, Cross, Tabernacle, Altar Card, and Veil Clarifications (1970/1975)}{Crux, tabernaculum, tabellae et velum}{1970--2002/2008 --- CDWDS Response (25 Sept 2000, Prot. N. 2036/00/L); GIRM (2002/2008) nos. 118c, 299, 306, 308, 314--315}{N43}
The third typical edition of the Roman Missal (2002/2008) together with a formal Roman clarification in 2000 resolved four sanctuary ambiguities inherited from the 1970 and 1975 editions. First, whereas 1975 GIRM no. 270 spoke only of a cross visible to the assembly without specifying a corpus---producing widespread bare crosses---2002/2008 GIRM no. 308 explicitly mandates a cross bearing the figure of Christ crucified (cum figura Christi crucifixi) on or near the altar, remaining there appropriately even outside Mass. Second, when the freestanding altar wording of GIRM no. 299 raised questions about historic churches, the Congregation for Divine Worship and the Discipline of the Sacraments issued its Response of 25 September 2000 (Prot. N. 2036/00/L; Notitiae 36:171--173), clarifying that expedit expresses desirability rather than an absolute ban on historic wall-attached altars or celebration ad orientem. Third, 2002/2008 GIRM nos. 314--315 refined 1975 GIRM no. 276 by placing two tabernacle locations on equal footing under the judgment of the diocesan bishop: either in the sanctuary apart from the altar of celebration (including on a former altar no longer used for Mass) or in an organically connected, readily noticeable adoration chapel. Fourth, the three 1962 altar cards (tabellae secretarum) were omitted from the 1970 Missal and remain excluded from the mensa under GIRM no. 306, whereas 2002/2008 GIRM no. 118c restored explicit praise of covering the prepared chalice with a veil (color of the day or white), leaving the burse unnamed as a customary accessory.

\PCDField{Current discipline (2008/2011)}{The altar or processional cross must display the figure of Christ crucified (GIRM no. 308); freestanding altars are recommended where possible without banning historic wall altars or celebration ad orientem (GIRM no. 299; CDWDS 25 Sept 2000); the diocesan bishop judges between placing the tabernacle in the sanctuary apart from the altar of celebration or in an organically connected adoration chapel (GIRM nos. 314--315); altar cards are excluded from the mensa (GIRM no. 306); and covering the prepared chalice with a veil of the color of the day or white is explicitly praised, with the burse remaining an unnamed customary accessory (GIRM no. 118c).}
\PCDField{1962 comparison}{\PCDTag{materially changed} The 1962 Roman Missal required a central altar crucifix, ordinarily placed the tabernacle in the middle of the high altar, prescribed three altar cards during Mass, and mandated a silk chalice veil with matching burse. Postconciliar clarifications in 2000--2008 reaffirm the crucifix, permit freestanding or historic wall altars, place sanctuary and chapel tabernacle sites on equal footing under the diocesan bishop's judgment, supersede the three altar cards on the mensa, and commend the chalice veil as a praiseworthy optional practice while leaving the burse as an unnamed custom.}
\PCDEntryRule
